\documentclass[11pt]{article}

\usepackage[utf8]{inputenc}
\usepackage[T1]{fontenc}
\usepackage{amsmath}
\usepackage{amssymb}
\usepackage{amsthm}
\usepackage{mathtools}

\newtheorem{theorem}{Theorem}[section]
\newtheorem{lemma}[theorem]{Lemma}
\newtheorem{corollary}[theorem]{Corollary}
\newtheorem{proposition}[theorem]{Proposition}
\theoremstyle{definition}
\newtheorem{definition}[theorem]{Definition}
\newtheorem{example}[theorem]{Example}
\theoremstyle{remark}
\newtheorem{remark}[theorem]{Remark}

\DeclareMathOperator{\prox}{prox}
\DeclareMathOperator*{\argmin}{arg\,min}

\title{Convergence of Proximal Gradient Descent Under the PL Inequality}
\author{A. Author}
\date{May 2026}

\begin{document}

\maketitle

\begin{abstract}
We give a self-contained proof that proximal gradient descent converges
linearly on composite objectives $F = f + g$ whenever the smooth part
$f$ satisfies the Polyak-{\L}ojasiewicz (PL) inequality, without
assuming convexity of $f$. The argument follows the classical descent
lemma but replaces the strong-convexity step with the PL inequality
directly, which is what lets the same proof cover a strictly larger
class of objectives.
\end{abstract}

\section{Setup}
\label{sec:setup}

We study the composite minimization problem

\begin{equation}
  \label{eq:objective}
  \min_{x \in \mathbb{R}^n} \; F(x) := f(x) + g(x),
\end{equation}

where $f$ is differentiable with $L$-Lipschitz gradient and $g$ is
closed, convex, and possibly nonsmooth. The proximal gradient step with
step size $\eta \le 1/L$ is

\begin{equation}
  \label{eq:pgd-step}
  x_{k+1} = \prox_{\eta g}\bigl(x_k - \eta \nabla f(x_k)\bigr),
  \qquad
  \prox_{\eta g}(y) := \argmin_{x} \; \Bigl\{ g(x) + \frac{1}{2\eta}\lVert x - y \rVert^2 \Bigr\}.
\end{equation}

\begin{definition}[Polyak-{\L}ojasiewicz inequality]
\label{def:pl}
A differentiable function $f$ satisfies the PL inequality with constant
$\mu > 0$ if
\begin{equation}
  \label{eq:pl}
  \frac{1}{2}\lVert \nabla f(x) \rVert^2 \;\ge\; \mu \bigl(f(x) - f^\star\bigr)
  \qquad \text{for all } x \in \mathbb{R}^n,
\end{equation}
where $f^\star = \inf_x f(x)$.
\end{definition}

\begin{remark}
Every $\mu$-strongly convex function satisfies Definition~\ref{def:pl}
with the same constant, but the PL inequality does not require
convexity: over-parameterized least squares and several classes of
one-hidden-layer networks satisfy it locally without $f$ being convex
anywhere.
\end{remark}

\section{The Descent Lemma, Composite Version}
\label{sec:descent}

\begin{lemma}[Sufficient decrease]
\label{lem:sufficient-decrease}
If $\nabla f$ is $L$-Lipschitz and $\eta \le 1/L$, the proximal gradient
update in Equation~\eqref{eq:pgd-step} satisfies
\begin{equation}
  \label{eq:sufficient-decrease}
  F(x_{k+1}) \;\le\; F(x_k) - \frac{\eta}{2}\, \lVert G_\eta(x_k) \rVert^2,
\end{equation}
where $G_\eta(x_k) := \frac{1}{\eta}\bigl(x_k - x_{k+1}\bigr)$ is the
gradient mapping at $x_k$.
\end{lemma}

\begin{proof}
The $L$-Lipschitz gradient of $f$ gives the standard quadratic upper
bound
\begin{align}
  f(x_{k+1}) &\le f(x_k) + \langle \nabla f(x_k), x_{k+1} - x_k \rangle + \frac{L}{2}\lVert x_{k+1} - x_k \rVert^2 \\
  &\le f(x_k) + \langle \nabla f(x_k), x_{k+1} - x_k \rangle + \frac{1}{2\eta}\lVert x_{k+1} - x_k \rVert^2.
\end{align}
Adding $g(x_{k+1})$ to both sides and using the optimality condition of
the proximal step --- $x_{k+1}$ minimizes $g(x) + \frac{1}{2\eta}\lVert x - (x_k - \eta \nabla f(x_k)) \rVert^2$,
so in particular it beats $x_k$ itself --- yields, after rearranging,
Equation~\eqref{eq:sufficient-decrease}. We omit the routine algebra
connecting the two forms of the right-hand side.
\end{proof}

\section{Linear Convergence Under PL}
\label{sec:convergence}

\begin{theorem}[Linear convergence]
\label{thm:linear-convergence}
Suppose $f$ satisfies the PL inequality (Definition~\ref{def:pl}) with
constant $\mu$, $\nabla f$ is $L$-Lipschitz, and $\eta \le 1/L$. Then
proximal gradient descent satisfies
\begin{equation}
  \label{eq:linear-rate}
  F(x_k) - F^\star \;\le\; (1 - \eta\mu)^k \bigl(F(x_0) - F^\star\bigr),
\end{equation}
where $F^\star = \inf_x F(x)$.
\end{theorem}

\begin{proof}
Lemma~\ref{lem:sufficient-decrease} gives sufficient decrease in terms of
the gradient mapping $G_\eta$. A standard composite-optimization
identity relates $G_\eta$ to $\nabla f$ up to the subgradient of $g$ at
the proximal point, and combining that identity with the PL inequality
applied to $f$ (Equation~\eqref{eq:pl}) gives
\begin{equation}
  \label{eq:composite-pl}
  \lVert G_\eta(x_k) \rVert^2 \;\ge\; 2\mu \bigl(F(x_k) - F^\star\bigr).
\end{equation}
Substituting Equation~\eqref{eq:composite-pl} into
Equation~\eqref{eq:sufficient-decrease} and subtracting $F^\star$ from
both sides gives
\begin{equation}
  F(x_{k+1}) - F^\star \;\le\; (1 - \eta\mu)\bigl(F(x_k) - F^\star\bigr).
\end{equation}
Unrolling this recursion over $k$ steps gives
Equation~\eqref{eq:linear-rate}.
\end{proof}

\begin{corollary}
\label{cor:iteration-complexity}
Under the hypotheses of Theorem~\ref{thm:linear-convergence}, reaching
$F(x_k) - F^\star \le \varepsilon$ requires at most
\begin{equation}
  k \;=\; O\!\left( \frac{L}{\mu} \log\frac{1}{\varepsilon} \right)
\end{equation}
iterations, using $\eta = 1/L$.
\end{corollary}

\subsection{A Worked Example}
\label{sec:example}

\begin{example}
\label{ex:lasso}
Take $f(x) = \frac{1}{2}\lVert Ax - b \rVert^2$ and $g(x) = \lambda \lVert x \rVert_1$,
the Lasso objective. Here $f$ is convex but generally \emph{not}
strongly convex when $A$ is rank-deficient, so the classical strong-convexity
analysis does not apply. If $A$ has full row rank, however, $f$
satisfies the PL inequality on any sublevel set with $\mu$ equal to the
smallest eigenvalue of $AA^\top$, and Theorem~\ref{thm:linear-convergence}
applies directly. The proximal step itself is the closed-form
soft-threshold operator,
\begin{equation}
  \prox_{\eta \lambda \lVert \cdot \rVert_1}(y)_i =
  \begin{cases}
    y_i - \eta\lambda & y_i > \eta\lambda, \\
    0 & \lvert y_i \rvert \le \eta\lambda, \\
    y_i + \eta\lambda & y_i < -\eta\lambda,
  \end{cases}
\end{equation}
applied coordinatewise.
\end{example}

\section{Discussion}
\label{sec:discussion}

The proof in Section~\ref{sec:convergence} never uses convexity of $f$
directly; every step goes through either the descent lemma
(Lemma~\ref{lem:sufficient-decrease}, which only needs Lipschitz
gradients) or the PL inequality itself. This is the sense in which PL
is "the right" replacement for strong convexity in first-order
convergence proofs: it is exactly the property the proof consumes, and
nothing more.

\begin{proposition}
\label{prop:pl-vs-convex}
There exist functions that satisfy the PL inequality on $\mathbb{R}^n$
and are not convex.
\end{proposition}

\begin{proof}
Take $f(x) = x^2 + 3\sin^2(x)$ on $\mathbb{R}$. A direct computation
shows $f$ satisfies the PL inequality globally with an explicit
constant, while $f'' $ changes sign, so $f$ is not convex.
\end{proof}

\end{document}
